% Paper G: Structural Limits of Unification
% Layer III: Accessibility, Incompleteness, and a Programme-Level Final-Cut Diagnosis
% COMPLETE VERSION — Final Draft

\documentclass[12pt,a4paper]{article}
\usepackage{amsmath,amssymb,amsfonts,amsthm}
\usepackage{physics}
\usepackage{hyperref}
\usepackage{geometry}
\usepackage{array}
\newcolumntype{P}[1]{>{\raggedright\arraybackslash}p{#1}}
\usepackage{booktabs}
\geometry{margin=1in}

\newtheorem{theorem}{Theorem}[section]
\newtheorem{proposition}[theorem]{Proposition}
\newtheorem{lemma}[theorem]{Lemma}
\newtheorem{definition}[theorem]{Definition}
\newtheorem{remark}[theorem]{Remark}

\title{Structural Limits of Unification:\\
Accessibility, Incompleteness, and a Programme-Level Final-Cut Diagnosis}

\author{
  Sidong Liu, PhD \\
  iBioStratix Ltd \\
  \texttt{sidongliu@hotmail.com}
}

\date{July 2026}

\begin{document}
\emergencystretch=2em
\raggedbottom
\hbadness=5000
\vbadness=5000

\maketitle

\begin{abstract}
This paper offers a programme-level philosophy-of-science diagnosis of where an accessibility-based framework may encounter limits when asked to represent its own selection rules. The former structural non-self-grounding lemma is withdrawn: no typed fixed-point domain or map has been supplied, and partial internalization or nontrivial self-reference remains open. Regress and explanatory collapse are therefore illustrative scenarios under a particular explanatory demand, not an exhaustive theorem or a forced stopping point. The contribution is methodological: to identify the object/meta interface that a genuine self-grounding analysis would have to formalize and to motivate an architectural final cut without claiming its necessity.
\end{abstract}

\section{Introduction}
\label{sec:intro}

\subsection{The Expectation of Completeness}

The aspiration toward a unified physical description has historically been guided by the expectation that deeper unification corresponds to increased completeness. In this traditional view, apparent multiplicity---of forces, degrees of freedom, or explanatory principles---is taken to signal provisional fragmentation, to be resolved by a more fundamental theory. A Theory of Everything, in its strongest formulation, is therefore often assumed to be both unifying and self-grounding: it should not only subsume all known interactions under a single framework, but also account for the conditions under which its own descriptions are possible.

The present work examines a more restricted possibility: a given explanatory strategy may encounter a boundary when asked to represent its own selection rules.  The discussion is a philosophy-of-science analysis in the metalanguage of the HAFF programme, not a theorem about all formal systems or a new physical no-go result.  Its motivation comes from the preceding papers \cite{Liu2026PaperA,Liu2026PaperB,Liu2026PaperC,Liu2026PaperD,Liu2026PaperE,Liu2026PaperF}, which build on the non-uniqueness of tensor factorizations \cite{Zanardi2001,Zanardi2004}.

\subsection{Summary of the Preceding Framework}

Across Papers D--F, gravitational reconstruction, measurement records, and temporal evidence are discussed using related selection and stability structures.  The revised papers use differently typed objects and make differently graded claims; they do not jointly prove that the three phenomena emerge from one ``accessible algebra.''

\begin{itemize}
    \item \textbf{Paper D}: Geometric conclusions are conditional on modular and holographic reconstruction inputs.
    \item \textbf{Paper E}: Probe-reachable effects, pointer/record algebras, and exact dynamical fixed points are separated rather than intersected into one algebra.
    \item \textbf{Paper F}: Redundant records can grow in a supplied branching dynamics from a low-record boundary; the resulting order does not derive time from algebra alone.
\end{itemize}

These discussions share a programme-level question about how selection data enter physical description.  That common question is motivation, not a theorem-level unification.

\subsection{The Problem of Self-Grounding}

If the explanatory programme depends on restrictions---on probe reachability, record stability, modular invariance, or reconstruction dictionaries---then it is legitimate to ask whether those restrictions can themselves be represented inside the same formal object.  The answer depends on how object language, metalanguage, and encodings between them are specified; it is not fixed by the word ``accessibility.''

The aim of this final layer is therefore diagnostic: to articulate a possible regress when selection rules are repeatedly redescribed, to identify the extra formal data a genuine fixed-point theorem would require, and to explain why the present HAFF papers do not yet supply such a theorem.

\subsection{Scope and Limitations}

Several clarifications are necessary at the outset.

First, the incompleteness identified in this paper is not formal in the sense of G\"{o}del's incompleteness theorems. We do not claim that the framework contains undecidable propositions within a formal system, nor do we invoke metamathematical results. The incompleteness is \emph{architectural}: it concerns the explanatory roles within a physical framework, not provability within a formal calculus.

Second, we do not claim that the structural limits identified here apply universally to all conceivable approaches to unification. The present analysis is confined to the algebraic and coarse-graining-based framework developed in Papers A--F. Whether alternative formalisms---category-theoretic, non-algebraic, or radically background-free---would exhibit analogous limits is an open question that we do not address.

Third, the proposed stopping point is not temporal, existential, or normative.  It is a methodological option for the present programme.  A richer formalization could internalize some selection data without inconsistency, so necessity is not claimed.

\subsection{Structure of the Paper}

Section~\ref{sec:accessibility} distinguishes physical observables from syntax and meta-description. Section~\ref{sec:lemma} replaces the former no-go lemma with a heuristic non-self-grounding observation and states its missing hypotheses. Section~\ref{sec:collapse} gives an illustrative regress scenario rather than an exhaustive dichotomy. Section~\ref{sec:cut} presents the final cut as a methodological proposal. Section~\ref{sec:nonclaims} states explicit non-claims. Section~\ref{sec:conclusion} concludes with reflections on the proposal's significance.

\section{Accessibility as a Non-Global Constraint}
\label{sec:accessibility}

\subsection{The Role of Accessibility in HAFF}

The HAFF papers use several notions previously grouped under ``accessibility.''  Paper E now distinguishes probe-reachable effects, pointer/record sectors, and exact fixed points; Paper D uses selected modular-stable algebraic data; Paper F uses a declared fragmentation and record dynamics.  Treating all of these as one endofunction on a space of algebras would erase their types.

In Papers D--F, this point is developed across distinct domains:

\begin{itemize}
    \item Effective geometry is conditional on a reconstruction dictionary and modular hypotheses.
    \item Measurement objectivity concerns robust pointer records, not a proof of a unique outcome.
    \item Record directionality is conditional on dynamics, fragmentation, and a low-record boundary.
\end{itemize}

In each case, some selection data are supplied.  The types and epistemic status of those data must be tracked separately.

\subsection{Structural vs.\ Epistemic Constraints}

A crucial distinction must be drawn between epistemic and structural constraints. This distinction is contested in philosophy of physics, and we acknowledge that what follows adopts it as a working criterion rather than a demonstrated result.

\begin{definition}[Epistemic Constraint]
A constraint is \textbf{epistemic} if it concerns what can be known, inferred, or verified by agents, given their informational position.
\end{definition}

\begin{definition}[Structural Constraint]
A constraint is \textbf{structural} if it concerns what descriptions are well-defined, given a pattern of physical interactions, independent of any agent's knowledge or epistemic state.
\end{definition}

The HAFF programme intends these inputs to be structural when they are fixed by dynamics and states, but their independence from a chosen factorization, fragment family, probe class, or reconstruction dictionary remains a substantive mereological question.

Given all of those inputs, different agents can apply the same rule.  That conditional reproducibility does not by itself prove that the inputs are uniquely selected by the underlying physics.

\begin{remark}[Contested Distinction]
We acknowledge that this distinction between epistemic and structural constraints is philosophically contested. The framework does not claim to have resolved this broader debate.  The burden is constructive: a structural claim should exhibit the physical data and selection map that make the relevant output invariant under admissible redescriptions.
\end{remark}

\subsection{Non-Globality of Accessibility}

The present papers do not provide a privileged, all-encompassing selected algebra from which all effective descriptions are derived.  Each construction uses restrictions whose source is only partly analyzed.  This is an open interface problem, not yet a proof that internal specification is impossible.

Before discussing internalization, one type distinction is essential.  An element $O$ of an observable algebra represents a physical observable.  A formula, algorithm, or data structure specifying which algebra is selected belongs to a syntax or metalanguage until an explicit encoding maps it into physical degrees of freedom.  Algebra membership alone does not mean that an algebra ``contains its own description.''

\section{A Heuristic Observation on Self-Grounding}
\label{sec:lemma}

\subsection{Status of the Former Lemma}
\label{lem:nonsg}

The former ``Structural Non-Self-Grounding Lemma'' is withdrawn as a theorem.  No formal domain of descriptions, accessibility endomap, order structure, or continuity/monotonicity hypothesis was defined, so the claimed absence of nontrivial fixed points did not follow.

The surviving heuristic is narrower: if every explanation of a selection rule is required to be given by another rule of exactly the same, still-unformalized kind, a regress can be generated by repeatedly asking what selects the next rule.  This is a statement about a chosen explanatory demand.  It does not establish inconsistency, collapse, or the nonexistence of a self-representing model.

\subsection{Why the Fixed-Point Argument Does Not Prove a No-Go Result}

The earlier argument can be separated into the following typed questions.

\paragraph{Object level.}
An observable algebra contains operators and their limits.  It does not automatically contain formulas naming that algebra, proofs about its closure, or code implementing a selection rule.  Identifying an observable $\hat A$ with a syntactic description of ``accessibility'' is therefore a type error unless an encoding and decoding map are supplied.

\paragraph{Meta level.}
Selection criteria are presently written in the theorist's metalanguage.  One may ask whether a physical subsystem can encode them, but this requires a representation of syntax, an evaluation map, and a criterion for faithful self-application.  None is defined in HAFF Papers A--F.

\paragraph{Fixed points.}
Even after such a formalization, self-application does not imply nonexistence.  Fixed-point and diagonal results routinely produce nontrivial self-reference under suitable hypotheses: the Kleene recursion theorem, Lawvere's fixed-point formulation, Scott's domain-theoretic fixed points, and Tarski's lattice theorem are standard counterweights to a blanket exclusion \cite{Kleene1952,Lawvere1969,Scott1972,Tarski1955}.  They do not prove that HAFF has the desired fixed point; they show that its absence must be demonstrated from a specified map and domain rather than inferred from the word ``self-reference.''

\paragraph{What a genuine theorem would require.}
A mathematical non-self-grounding theorem would at least have to specify:
\begin{itemize}
    \item a space or category of descriptions and the type of each object;
    \item a selection or evaluation map with a well-defined codomain;
    \item the encoding that connects syntax or meta-data to physical observables;
    \item the admissible fixed points and why known fixed-point hypotheses fail; and
    \item a proved exhaustiveness statement if regress and totalization are claimed to be the only outcomes.
\end{itemize}

\begin{remark}[Heuristic Status]
Until those data are supplied, non-self-grounding is a plausible philosophy-of-science diagnosis of the programme's current explanatory architecture.  It is not a lemma, and alternatives such as partial internalization, stratified languages, reflective encodings, or nontrivial fixed points remain open.
\end{remark}

\section{An Illustrative Regress Scenario}
\label{sec:collapse}

To make the heuristic concrete, consider one particular extension that tries to internalize a selection rule without first formalizing the object/meta interface.  The example illustrates how a regress can arise; it does not prove that every internalization fails.

\subsection{Hypothetical Extension}

\paragraph{Step 1: Choose an encoding.}
Suppose a physical register $R$ is introduced to encode descriptions of candidate subalgebras and their selection scores.  The register observables are ordinary algebra elements; a separate map must say which register states encode which syntactic descriptions.

\paragraph{Step 2: Supply dynamics and evaluation.}
For the encoding to do explanatory work, one must specify:
\begin{itemize}
    \item dynamics for $R$ and its coupling to the system; and
    \item an evaluation procedure that maps an encoded rule and physical data to a selected output.
\end{itemize}

\paragraph{Step 3: Reapply the same explanatory demand.}
If the programme insists that the encoding, evaluation map, and register dynamics must themselves be selected by another rule of the same kind, then a new datum is introduced at each level.  Repeating that demand produces a regress by construction.

\subsection{Two Possible Outcomes, Not an Exhaustive Dichotomy}

\paragraph{Outcome (a): Regress under unrestricted reapplication.}
If no base rule, fixed point, equivalence class, or stopping convention is admitted, the explanatory request repeats indefinitely.  This is a consequence of the stipulated demand, not a contradiction in the observable algebra.

\paragraph{Outcome (b): Failure of a naive totalization.}
One attempted stopping convention is to declare every effect or observable equally selected.  In models where the physical content depends on a nontrivial probe class, pointer sector, or reconstruction subalgebra, that convention erases the contrast being modeled:
\begin{itemize}
    \item no nontrivial selection remains;
    \item pointer records are no longer distinguished from arbitrary observables; and
    \item any geometry or record order conditional on the selected structure loses that input.
\end{itemize}
This shows that one totalization proposal is unhelpful.  It does not exclude typed fixed points, reflective sublanguages, externally supplied base data, or other ways of terminating the regress.

\subsection{Conclusion of the Scenario}

The scenario supports a methodological warning: internalization must preserve the object/meta type distinction and specify its evaluation map.  Without those ingredients, claims of either self-grounding or necessary collapse are premature.

\section{A Proposed Final Cut}
\label{sec:cut}

\subsection{Stopping as a Methodological Choice}

The preceding analysis motivates a limited sense in which the current HAFF programme is incomplete: it has not formalized how all of its selection inputs are themselves chosen.  This is neither a formal incompleteness theorem nor a metaphysical limit, and it may be addressed by future interface constructions.

Within the programme, geometry, measurement records, and temporal records all use supplied selection or stability data, but of different types and at different evidential grades.  The present paper asks where to stop redescribing those inputs; it does not prove that no further physical account can be given.

\subsection{The Final Cut}

The notion of a ``final cut'' is introduced as a proposed bookkeeping boundary. It does not denote a temporal endpoint or the completion of physics.  It marks where the current papers stop and expose their selection data as inputs rather than presenting a further unproved closure claim.

This cut is revisable.  A future typed model of self-description, a nontrivial fixed point, or a physical derivation of the selection data could move it.

Beyond this point, a further extension must state how descriptions of selection rules are represented and evaluated.  The heuristic in Section~\ref{sec:lemma} identifies that obligation but does not show that it admits no nondegenerate solution.

The stopping point is therefore selected provisionally so that an open interface is not disguised as a theorem.

\subsection{Incompleteness as Internal Limit}

In this sense, the present incompleteness is programme-relative and provisional: the formalism has not supplied the typed self-description data listed above.

The phrase ``unselected structure'' is therefore interpretive shorthand, not a demonstrated domain that admits no geometry, temporal ordering, or observational standpoint.

A theory can be clearer by marking where its current derivations stop and its interface assumptions begin.

\begin{remark}[Open Alternatives]
Section~\ref{sec:collapse} gives two outcomes for one naive internalization strategy, not an exhaustive dichotomy.  A proposed extension should specify how it handles encoding, evaluation, and repeated selection; whether it does so by a base datum, stratification, or fixed point is left open.
\end{remark}

\section{Scope and Non-Claims}
\label{sec:nonclaims}

To prevent misinterpretation, several non-claims must be stated explicitly.

\begin{enumerate}
    \item \textbf{No invocation of G\"{o}del's theorems.} This work does not invoke G\"{o}del's incompleteness theorems, nor does it rely on any formal analogy to them. The incompleteness identified here is architectural, not metamathematical. It concerns explanatory roles within a physical framework, not provability within a formal system.
    
    \item \textbf{No claim of universal applicability.} No claim is made that accessibility constraints apply universally across all conceivable approaches to unification. The present analysis is confined to the algebraic and coarse-graining-based framework developed in Papers A--F.
    
    \item \textbf{Contested distinction acknowledged.} The distinction between structural and epistemic constraints is acknowledged to be contested in the philosophy of physics. The framework adopts this distinction as a working criterion, grounded in interaction-determined stability conditions. It does not claim to have resolved the broader philosophical debate.
    
    \item \textbf{No foreclosure of further formalization.} The proposed final cut does not preclude further physical progress or a typed self-description model within HAFF.  Further internal extension has not been shown incoherent.
    
    \item \textbf{No uniqueness or necessity claim.} Other frameworks may stop elsewhere, and future versions of this programme may move or remove the proposed cut.
    
    \item \textbf{No philosophical finality.} This is a structural observation within a specific formalism, not a philosophical conclusion about the ultimate nature of reality or the limits of knowledge as such. The stopping point identified is methodological and structural, not existential.
    
    \item \textbf{No consciousness or observer-creation claims.} The framework does not claim that consciousness plays a fundamental role, that observers create reality, or that accessibility is observer-relative in any subjective sense.
    
    \item \textbf{No claim of novelty regarding theoretical presupposition.} We acknowledge that explanatory frameworks can have presuppositions they do not presently explain.  The contribution here is to identify where this issue appears in HAFF and what a formal replacement for the withdrawn lemma would need.
\end{enumerate}

Any downstream ``No View from Nowhere'' closure principle that relied on the former lemma must therefore be conditionalized.  Likewise, any later Paper~III use of this discussion is motivation-only unless it supplies the missing typed fixed-point argument.  No downstream theorem is established here.

\section{Conclusion}
\label{sec:conclusion}

\subsection{Summary of Results}

This paper has examined conditions under which the HAFF framework may encounter a boundary in extending its own explanatory strategy.

The central results are:

\begin{enumerate}
    \item \textbf{Accessibility-specification asymmetry}: HAFF descriptions use selection and stability data that are presently stated in a theorist's metalanguage rather than derived from one common object-level formalism.
    
    \item \textbf{Structural non-self-grounding (programme-level diagnosis)}: The current framework exhibits an asymmetry between accessibility-specification and accessibility-governed descriptions; a formal fixed-point exclusion remains open (Section~\ref{sec:lemma}).
    
    \item \textbf{Illustrative regress scenario}: Repeatedly demanding a selector of the same unformalized kind can generate regress or totalization pressure. This does not show that every typed internalization fails.
    
    \item \textbf{Architectural hypothesis}: A final cut may serve as a programme-level stopping rule while the formal self-grounding problem remains unresolved; it is structural motivation, not a proved necessity.
\end{enumerate}

\subsection{Methodological Significance}

Unification is often equated with the elimination of boundaries. The analysis presented here suggests a different criterion: that a unifying framework should distinguish between boundaries that are provisional and those that are structural.

Papers D--F motivate related boundary questions in the treatment of gravity, measurement, and temporal records. This final layer proposes that making those open interfaces explicit can improve explanatory clarity without turning them into an impossibility theorem.

The value of the framework lies not in claiming to determine what cannot be said, but in identifying which selection and self-description interfaces remain unformalized. Treating a final cut as an architectural hypothesis keeps that limitation visible while leaving stronger fixed-point constructions open.

\subsection{Open Questions}

Several questions remain beyond the scope of this analysis:

\begin{itemize}
    \item Whether alternative frameworks (category-theoretic, non-algebraic, or background-free) would exhibit analogous structural limits.
    \item Whether the structural/epistemic distinction adopted here can be given a more robust philosophical foundation.
    \item Whether the collapse scenario admits any non-trivial avoidance strategies not considered here.
\end{itemize}

These questions are left for future investigation.

\section*{Acknowledgments}

The author thanks the anonymous reviewers for their insightful comments and suggestions, which greatly improved the clarity and rigor of this work.

\begin{thebibliography}{99}

\bibitem{Liu2026PaperA}
S.~Liu, \emph{Emergent Geometry from Coarse-Grained Observable Algebras: The Holographic Alaya-Field Framework}, Zenodo (2026), DOI: 10.5281/zenodo.18361706.

\bibitem{Liu2026PaperB}
S.~Liu, \emph{Accessibility, Stability, and Emergent Geometry: Conceptual Clarifications on the Holographic Alaya-Field Framework}, Zenodo (2026), DOI: 10.5281/zenodo.18367060.

\bibitem{Liu2026PaperC}
S.~Liu, \emph{Causation, Agency, and Existence: Structural Constraints and Interpretive Bridges}, Zenodo (2026), DOI: 10.5281/zenodo.18374805.

\bibitem{Liu2026PaperD}
S.~Liu, \emph{Observable-Algebra Selection and Emergent Gravitational Phenomena: A Conditional Structural Analysis}, Zenodo (2026), DOI: 10.5281/zenodo.18388881.

\bibitem{Liu2026PaperE}
S.~Liu, \emph{Measurement as Accessibility: A Structural Analysis of Observable Algebra Selection}, Zenodo (2026), DOI: 10.5281/zenodo.18400065.

\bibitem{Liu2026PaperF}
S.~Liu, \emph{Accessibility Propagation and Temporal Direction: A Conditional Structural Analysis}, Zenodo (2026), DOI: 10.5281/zenodo.18400425.

\bibitem{Zanardi2001}
P. Zanardi, \emph{Virtual Quantum Subsystems}, Phys. Rev. Lett. \textbf{87}, 077901 (2001).

\bibitem{Zanardi2004}
P. Zanardi, D. A. Lidar, and S. Lloyd, \emph{Quantum Tensor Product Structures are Observable Induced}, Phys. Rev. Lett. \textbf{92}, 060402 (2004).

\bibitem{Kleene1952}
S. C. Kleene, \emph{Introduction to Metamathematics}, North-Holland (1952).

\bibitem{Lawvere1969}
F. W. Lawvere, \emph{Diagonal arguments and cartesian closed categories}, in \emph{Category Theory, Homology Theory and their Applications II}, Lecture Notes in Mathematics \textbf{92}, 134--145 (1969).

\bibitem{Scott1972}
D. Scott, \emph{Continuous lattices}, in \emph{Toposes, Algebraic Geometry and Logic}, Lecture Notes in Mathematics \textbf{274}, 97--136 (1972).

\bibitem{Tarski1955}
A. Tarski, \emph{A lattice-theoretical fixpoint theorem and its applications}, Pacific J. Math. \textbf{5}, 285--309 (1955).

\end{thebibliography}

\end{document}
